\section{Discussion}
\label{sec:discussion}

In this section,
we discuss the key aspects and implications of our work.

\paragraph{Applicability to other \SE engines}
% our work provide the tools to address semantic gaps
% for performance, depends on the gap size
While our current implementation targets \klee, %our underlying architecture is modular.
we believe that our approach can be applied to other \SE engines whose semantics differ from those of \klee.
\klee's symbolic semantics significantly differs from our symbolic semantics (\Cref{sec:symbolic-semantics}),
and our approach provides the mechanisms to bridge such gaps.
First, we develop a theory that allows us to prove that the SMT transformations in \klee preserve semantic equivalence.
Second, we use interface lemmas which help to bridge the gap between \klee's dual representation of memory objects and our symbolic semantics.

%Regarding the impact on performance (proof compilation time, \etc),
%In general, if there is a substantial semantic gap,
%the manual proof effort and the proof checking cost are likely to increase.
%Otherwise, we expect the performance to be similar.

\paragraph{...}

\paragraph{...}

% TODO: applying our approach to other SE tools / semantics
% TODO: modeling of malloc
% TODO: the mechanized theory can by useful in other contexts
